\section{Motivation}

During Economy Week (Deutsch: Wirtschaftwoche), as an exercise in \textit{business administration}, we "played"
a game, where you were the board of a company manufacturing some arbitrary product and had to make decisions to
increase sales, lower costs and generally beat out your competition. For me, the aspect of being able to come up with
novel management ideas and compete with your friends really hooked me into the concept. This may sound hyperbolic, but
during that week I genuinely woke up excited to go to school, not to actually learn stuff (and whatever) but just to
play a few more rounds and (aspirationally) completely destroy my competition.

The game we got to play during Economy Week, in my opinion, had a tantalising hook,. but the interface was utterly
hostile to non-directed interaction: it was pretty much an Excel table! I set off to build a better, deeper
version of this simulation.

\begin{quote}
	“We do this not because it is easy, but because we thought it would be easy” \\
	\null \hfill \textbf{Not JFK.}
\end{quote}

This is the running motto of this project.

Reflecting the main intention of this project, the name of the game is (currently) \textbf{Wisim} (German:
\textbf{Wir}tschafts \textbf{Sim}ulation). From this point onward, all mentions of "Wisim" refer to the product of this
project.
